\chapter{数组、双指针和滑动窗口}

数组题是刷题的地基。双指针和滑动窗口不是孤立技巧，而是从暴力枚举子结构中优化出来的方法。双指针常用于有序数组、两端收缩和左右边界协作；滑动窗口常用于连续子数组或子串，并且窗口状态能够增量维护。

盛最多水的容器暴力解枚举左右边界。面积由宽度和较短边决定，所以每个下标对都可能成为候选。

\begin{lstlisting}[language=C++,caption={盛水容器暴力解}]
int max_area_bruteforce(const std::vector<int>& height)
{
    int best = 0;
    for (int left = 0; left < static_cast<int>(height.size()); ++left) {
        for (int right = left + 1; right < static_cast<int>(height.size()); ++right) {
            const int width = right - left;
            const int bounded_height = std::min(height[left], height[right]);
            best = std::max(best, width * bounded_height);
        }
    }
    return best;
}
\end{lstlisting}

暴力法是 \complexity{O(n^2)}。优化来自候选淘汰。初始取最宽的左右边界，如果移动较高的一边，宽度变小，而较短边没有变高，面积不可能变大。因此只能移动较短边，才有机会提高瓶颈高度。

\begin{lstlisting}[language=C++,caption={盛水容器双指针解}]
int max_area_two_pointers(const std::vector<int>& height)
{
    int left = 0;
    int right = static_cast<int>(height.size()) - 1;
    int best = 0;

    while (left < right) {
        const int width = right - left;
        const int bounded_height = std::min(height[left], height[right]);
        best = std::max(best, width * bounded_height);

        if (height[left] <= height[right]) {
            ++left;
        } else {
            --right;
        }
    }
    return best;
}
\end{lstlisting}

不变量是：每一步排除的短边，不会和更内侧的任何边形成更优答案，因为宽度只会变小，短边高度仍然是瓶颈。这个证明比“左右指针往中间走”重要得多。

长度最小的子数组展示滑动窗口。暴力法枚举所有起点和终点，遇到和达到目标时更新答案。它慢在每个起点都重新累加。

\begin{lstlisting}[language=C++,caption={长度最小子数组滑动窗口}]
int min_subarray_len_sliding_window(int target, const std::vector<int>& nums)
{
    int best = std::numeric_limits<int>::max();
    int left = 0;
    int sum = 0;

    for (int right = 0; right < static_cast<int>(nums.size()); ++right) {
        sum += nums[right];
        while (sum >= target) {
            best = std::min(best, right - left + 1);
            sum -= nums[left];
            ++left;
        }
    }

    return best == std::numeric_limits<int>::max() ? 0 : best;
}
\end{lstlisting}

这个优化依赖正整数。右边界右移时窗口和只会增加，左边界右移时窗口和只会减少。如果数组包含负数，这个单调性消失，滑动窗口可能失效，需要考虑前缀和、单调队列或其他结构。

无重复字符的最长子串也是窗口题，但窗口状态不是和，而是字符最后出现位置。窗口不变量是 \code{[left, right]} 内没有重复字符。

\begin{lstlisting}[language=C++,caption={无重复字符最长子串}]
int longest_substring_without_repeat_window(const std::string& s)
{
    constexpr int ASCII_RANGE = 256;
    std::array<int, ASCII_RANGE> last_seen{};
    last_seen.fill(-1);

    int left = 0;
    int best = 0;

    for (int right = 0; right < static_cast<int>(s.size()); ++right) {
        const auto ch = static_cast<unsigned char>(s[right]);
        if (last_seen[ch] >= left) {
            left = last_seen[ch] + 1;
        }
        last_seen[ch] = right;
        best = std::max(best, right - left + 1);
    }

    return best;
}
\end{lstlisting}

这里用 \code{std::array<int, 256>} 是因为按 ASCII 字符处理。如果题目要求完整 Unicode 字符，需要先明确字符单元，或使用哈希表保存字符到位置。

三数之和是双指针的第二个标准模型。题目要求找出所有不重复三元组，使得三个数之和为 0。暴力方法枚举三个下标，时间复杂度 \complexity{O(n^3)}，还要处理去重；如果用哈希表固定两数再找第三数，可以降到 \complexity{O(n^2)}，但去重和重复值处理仍然复杂。真正稳定的做法是先排序，再固定第一个数，把剩下两个数变成有序数组上的两数之和。

排序后的关键性质是：固定 \code{first} 后，如果 \code{nums[left] + nums[right]} 太小，说明左端值不够大，只能右移 \code{left}；如果和太大，说明右端值太大，只能左移 \code{right}。这就是有序性带来的候选淘汰。去重也依赖排序：同一层循环里，如果当前值和前一个值相同，就跳过；找到一个答案后，左右指针都要跳过连续重复值。

\begin{lstlisting}[language=C++,caption={三数之和：排序加双指针}]
std::vector<std::vector<int>> three_sum(std::vector<int> nums)
{
    std::sort(nums.begin(), nums.end());
    std::vector<std::vector<int>> answer;

    for (int first = 0; first < static_cast<int>(nums.size()); ++first) {
        if (first > 0 && nums[first] == nums[first - 1]) {
            continue;
        }
        if (nums[first] > 0) {
            break;
        }

        int left = first + 1;
        int right = static_cast<int>(nums.size()) - 1;
        while (left < right) {
            const int sum = nums[first] + nums[left] + nums[right];
            if (sum == 0) {
                answer.push_back({nums[first], nums[left], nums[right]});
                const int left_value = nums[left];
                const int right_value = nums[right];
                while (left < right && nums[left] == left_value) {
                    ++left;
                }
                while (left < right && nums[right] == right_value) {
                    --right;
                }
            } else if (sum < 0) {
                ++left;
            } else {
                --right;
            }
        }
    }

    return answer;
}
\end{lstlisting}

这里参数按值传入，是因为函数内部要排序。如果不想修改调用者的数组，按值传参是合理的；如果题目允许原地修改，也可以传引用。复杂度是排序 \complexity{O(n\log n)} 加双指针枚举 \complexity{O(n^2)}，总时间 \complexity{O(n^2)}，额外空间不计答案为 \complexity{O(1)} 或排序栈空间。面试表达要强调：去重不是最后把答案放进 set，而是在生成答案时就避免重复，这样更可控。

找到字符串中所有字母异位词是固定长度窗口。题目给定字符串 \code{s} 和模式串 \code{p}，要求找出 \code{s} 中所有长度等于 \code{p} 且字符频次完全相同的起点。暴力方法枚举每个长度为 \code{p.size()} 的子串，然后排序比较或统计比较。如果每次排序，复杂度是 \complexity{O(nk\log k)}；如果每次重新统计，复杂度是 \complexity{O(nk)}。慢点在于相邻窗口高度重叠，却重复计算了大部分字符频次。

优化方法是维护固定长度窗口的频次数组。窗口右端每进入一个字符，就把它的计数加一；当窗口长度超过模式串长度，左端字符离开，计数减一。由于题目通常限定小写英文，可以用 \code{std::array<int, 26>}，比较两个数组是常数时间。窗口不变量是：处理到 \code{right} 时，\code{window} 正好记录当前长度不超过模式串长度的后缀窗口字符频次。

\begin{lstlisting}[language=C++,caption={找到所有字母异位词：固定长度窗口}]
std::vector<int> find_anagrams_window(const std::string& s,
                                      const std::string& p)
{
    constexpr int ALPHABET_SIZE = 26;
    std::vector<int> answer;
    if (s.size() < p.size()) {
        return answer;
    }

    std::array<int, ALPHABET_SIZE> need{};
    std::array<int, ALPHABET_SIZE> window{};
    for (const char ch : p) {
        ++need[ch - 'a'];
    }

    for (std::size_t right = 0; right < s.size(); ++right) {
        ++window[s[right] - 'a'];
        if (right >= p.size()) {
            --window[s[right - p.size()] - 'a'];
        }
        if (window == need) {
            answer.push_back(static_cast<int>(right + 1 - p.size()));
        }
    }

    return answer;
}
\end{lstlisting}

这题容易错在下标。窗口第一次达到目标长度是在 \code{right + 1 == p.size()}，起点是 \code{right + 1 - p.size()}。如果用 \code{int} 写下标，要小心 \code{s.size()} 是无符号类型；本书示例这里使用 \code{std::size_t} 控制循环，最终起点再转成 \code{int}，因为力扣返回 \code{vector<int>}。

最小覆盖子串是变长窗口里最重要的一题。题目要求在 \code{s} 中找最短子串，覆盖 \code{t} 中所有字符及其次数。暴力方法枚举所有子串，再检查是否覆盖，复杂度至少 \complexity{O(n^2 |\Sigma|)}。优化的核心是：右边界扩张直到窗口满足覆盖条件，然后尽可能移动左边界缩短窗口；一旦左移导致窗口不满足，再继续右移。这个过程成立，因为右移只会增加字符，左移只会删除字符，覆盖关系可以被增量维护。

\begin{lstlisting}[language=C++,caption={最小覆盖子串：变长窗口}]
std::string min_window_substring(const std::string& s, const std::string& t)
{
    constexpr int ASCII_RANGE = 128;
    std::array<int, ASCII_RANGE> need{};
    std::array<int, ASCII_RANGE> window{};
    int required = 0;

    for (const char ch : t) {
        const auto index = static_cast<unsigned char>(ch);
        if (need[index] == 0) {
            ++required;
        }
        ++need[index];
    }

    int formed = 0;
    int left = 0;
    int best_left = 0;
    int best_length = std::numeric_limits<int>::max();

    for (int right = 0; right < static_cast<int>(s.size()); ++right) {
        const auto in = static_cast<unsigned char>(s[right]);
        ++window[in];
        if (need[in] > 0 && window[in] == need[in]) {
            ++formed;
        }

        while (formed == required) {
            const int length = right - left + 1;
            if (length < best_length) {
                best_length = length;
                best_left = left;
            }

            const auto out = static_cast<unsigned char>(s[left]);
            --window[out];
            if (need[out] > 0 && window[out] < need[out]) {
                --formed;
            }
            ++left;
        }
    }

    if (best_length == std::numeric_limits<int>::max()) {
        return "";
    }
    return s.substr(best_left, best_length);
}
\end{lstlisting}

\code{required} 表示需要满足的不同字符种类数，\code{formed} 表示当前窗口里已经满足需求的字符种类数。不能只用窗口总长度判断，因为 \code{t = "AABC"} 时，两个 \code{A} 都必须覆盖。复杂度是 \complexity{O(n + m + |\Sigma|)}，在 ASCII 数组固定时可视为 \complexity{O(n + m)}。面试表达要说清楚：这题不是固定窗口，因为答案长度未知；它是“先满足，再收缩”的最短覆盖模型。

\begin{principlebox}{本章题目复盘模板}
数组和窗口题复盘时至少写五句话：暴力枚举什么；慢点是否来自重复扫描、重复统计或候选过多；优化使用了有序性、窗口增量状态还是候选淘汰；左右指针移动条件是什么；每个元素为什么只会被处理常数次。只写“用双指针”或“用滑动窗口”不算会。
\end{principlebox}

本章力扣主线：11 盛最多水的容器、15 三数之和、209 长度最小的子数组、3 无重复字符的最长子串、438 找到字符串中所有字母异位词、76 最小覆盖子串。复盘时必须写出窗口不变量和指针移动条件。
